\chapter*{\textbf{Ucapan Terima Kasih}}
\addcontentsline{toc}{chapter}{\MakeUppercase{Ucapan Terima Kasih}}

Bagian ucapan terima kasih digunakan untuk menyampaikan penghargaan kepada pihak-pihak yang telah memberikan bantuan, dukungan, bimbingan, fasilitas, maupun kontribusi selama pelaksanaan penelitian dan penyusunan tugas akhir.

Ucapan terima kasih ditulis secara formal, singkat, dan proporsional. Pihak yang disebutkan umumnya meliputi orang tua dan keluarga, dosen pembimbing, dosen penguji, program studi atau laboratorium, mitra penelitian, serta pihak lain yang memiliki kontribusi langsung terhadap penyelesaian tugas akhir.

Dalam penulisannya, gunakan bahasa yang santun dan profesional serta hindari ungkapan yang berlebihan, bersifat informal, atau tidak berkaitan dengan pelaksanaan penelitian. Bagian ini tidak digunakan untuk membahas hasil penelitian maupun kesimpulan tugas akhir.
